\subsection{The intermediate value theorem}

Continuity has consequences that go beyond evaluating limits.

\begin{theorembox}[Intermediate value theorem]
If $f$ is continuous on $[a,b]$, then every value between $f(a)$ and $f(b)$ is attained at some point of the interval.
\end{theorembox}

\begin{remarkbox}[Proof status]
The geometric mechanism is that a continuous graph cannot jump over an intermediate level. A full proof uses completeness of $\mathbb R$, for example through nested intervals or a least-upper-bound argument, and is not developed here.
\end{remarkbox}

\begin{examplebox}[Proving that a root exists]
Let
\[
f(x)=x^3-x-1.
\]
The function is continuous everywhere. Moreover,
\[
f(1)=-1,
\qquad
f(2)=5.
\]
Since $0$ lies between $-1$ and $5$, there is at least one $c\in(1,2)$ such that
\[
f(c)=0.
\]
This argument proves existence without finding an exact formula for the root.
\end{examplebox}

\begin{interpretationbox}[Continuity prevents jumps over values]
The theorem turns a visible property of an unbroken graph into a precise existence statement. Its purpose is not to compute the solution but to guarantee that a solution must occur.
\end{interpretationbox}

\begin{remarkbox}
The theorem guarantees at least one solution, not uniqueness. Additional information, such as monotonicity, is needed to prove that the solution is unique.
\end{remarkbox}
